% !TeX root = ../main.tex
%===========================================================================%
%->> Appendix A
%===========================================================================%

\appendix
\chapter{系統級模擬完整參數設定 (System-Level Simulation Parameter Settings)}
\label{appendix:A}

本附錄詳細列出本研究所採用之系統級模擬（SLS）設定參數。為了確保模擬環境之產業通用性、避免標準術語在翻譯上的歧義，並與 3GPP TR 38.864 網路節能（NES）技術報告中之 Annex B Set 1（Macro Scenario）基準規範對齊，本參數表之技術指標與配置內容皆以國際通訊標準英文術語進行完整表列。因技術參數項目繁多，特將表格區分為兩部分呈現：表 \ref{tab:appendix_3gpp_set1_part1} 主要涵蓋網路佈署拓撲與基站實體層配置 (Part I)，表 \ref{tab:appendix_3gpp_set1_part2} 則涵蓋終端參數、資源調度與隨機流量模型 (Part II)。

%--- 表格第一部分：網路佈署與基地台參數 ---
\begin{table}[htbp]
\centering
\caption{Simulation Assumptions Aligned with 3GPP TR 38.864 Annex B Set 1 (Part I)}
\label{tab:appendix_3gpp_set1_part1}
\renewcommand{\arraystretch}{1.25}
\begin{tabularx}{\textwidth}{|p{0.38\textwidth}|X|}
\hline
\textbf{Parameter Category} & \textbf{Simulation Value / Configuration (Set 1)} \\ \hline
\multicolumn{2}{|c|}{\textbf{Network Layout, Topology, and Mobility}} \\ \hline
Channel model & 3D-UMa (Urban Macro), based on 3GPP TR 38.901 \\ \hline
Inter-site distance (ISD) & 500 m \\ \hline
Network topology & 7 sites, 3 sectors per site (21 cells total) \\ \hline
Boundary handling & 3D Wrap-around layout \\ \hline
Carrier frequency & 4.0 GHz \\ \hline
System bandwidth & 100 MHz \\ \hline
Subcarrier spacing (SCS) & 30 kHz (Numerology $\mu = 1$) \\ \hline
Frame structure & TDD mode, 1 slot = 0.5 ms (14 OFDM symbols/slot) \\ \hline
Building penetration loss & 100\% Low Loss Model \\ \hline
UE spatial deployment & 80\% Indoor, 20\% Outdoor random distribution \\ \hline
UE mobility & Speed = 3 km/h, random directions \\ \hline
\multicolumn{2}{|c|}{\textbf{BS Antenna Array and Physical Layer Configuration}} \\ \hline
Antenna configurations  at TRxP & For 64T:\newline $(M, N, P, M_g, N_g, M_p, N_p) = (8, 8, 2, 1, 1, 4, 8)$ \newline $(d_H, d_V) = (0.5, 0.8)\lambda$ \\ \hline
Antenna element radiation pattern & Directional macro pattern per TR 38.864 (max gain: 8 dBi) \\ \hline
Maximum downlink Tx power & 49 dBm (with dynamic power boosting) \\ \hline
BS antenna height & 25 m \\ \hline
BS receiver noise figure & 5 dB \\ \hline
\multicolumn{2}{|c|}{\textbf{User Equipment (UE) Parameters}} \\ \hline
UE antenna configuration & 
  For 4R: $(M, N, P, M_g, N_g, M_p, N_p) = (1, 2, 2, 1, 1, 1, 2)$  \newline $(d_H, d_V) = (0.5, N/A)\lambda$ \\ \hline

UE antenna radiation pattern & Omnidirectional pattern (0 dBi element gain) \\ \hline
UE height & 1.5 m \\ \hline
UE maximum Tx power & 23 dBm \\ \hline
UE noise figure & 9 dB \\ \hline
\end{tabularx}
\end{table}

\clearpage % 強制換頁

%--- 表格第二部分：終端、傳輸與流量模型 ---
\begin{table}[htbp]
\centering
\caption{Simulation Assumptions Aligned with 3GPP TR 38.864 Annex B Set 1 (Part II)}
\label{tab:appendix_3gpp_set1_part2}
\renewcommand{\arraystretch}{1.25}
\begin{tabularx}{\textwidth}{|p{0.38\textwidth}|X|}
\hline
\textbf{Parameter Category} & \textbf{Simulation Value / Configuration (Set 1)} \\ \hline
\multicolumn{2}{|c|}{\textbf{Transmission Control and CSI Management}} \\ \hline
DL Transmission scheme & SU-MIMO spatial multiplexing (up to 4 layers adaptive selection) \\ \hline
Maximum DL modulation order & Up to 256 QAM \\ \hline
DL Channel estimation & Precoded CSI-RS based \\ \hline
DL Codebook & non-PMI transmission \\ \hline
Rank adaptation scheme & Fully dynamic rank adaptation (\texttt{uFixDLRank = 0}) \\ \hline
CSI periodicity & $T_{\text{CSI}} = 5$ ms \\ \hline
SDA delay & 2 slots \\ \hline
DL Receiver architecture & MMSE-IRC \\ \hline
Resource scheduling algorithm & Proportional Fair (PF) scheduling \\ \hline
\multicolumn{2}{|c|}{\textbf{Traffic Model, Upper Layers, and Power Baseline}} \\ \hline
Traffic model & FTP Model 3 \\ \hline
Packet size & $0.5 Mbytes$ \\ \hline
Mean inter-arrival time & $200 ms$\\ \hline
Max HARQ retransmission & 3 \\ \hline
Target block error rate (BLER) & 10\% for initial transmission \\ \hline
BS power consumption baseline & 3GPP TR 38.864 Category 1 power model \\ \hline
\end{tabularx}
\end{table}